% =============================================================================
\section{Task Selection}
\label{sec:tasksel}
% =============================================================================

Choosing tasks that are both ecologically valid and shared between SLM and LLM
user communities is essential for meaningful capability comparison.
We conducted a formative Reddit study to ground task selection empirically
rather than relying on researcher intuition.

% -----------------------------------------------------------------------------
\subsection{Data Collection}
% -----------------------------------------------------------------------------

We scraped Reddit via the PullPush API over a twelve-month window
(March~2025--March~2026) using a two-stream keyword strategy designed to
capture distinct but potentially overlapping user populations.

\textbf{Stream~1} (SLM/local) targeted small-model discourse with 12 anchor
terms (e.g., \textit{local llm}, \textit{small language model},
\textit{edge ai}, \textit{self-hosted llm}) in the subreddits
\textit{r/LocalLLaMA}, \textit{r/LocalLLM}, and \textit{r/SelfHosted}.
\textbf{Stream~2} (LLM/general) targeted large-model discourse with 12 anchor
terms (e.g., \textit{chatgpt}, \textit{generative ai},
\textit{foundation model}) in \textit{r/ChatGPT}, \textit{r/OpenAI},
\textit{r/ClaudeAI}, \textit{r/LLM}, and \textit{r/Gemini}.
Both streams applied a post-retrieval substring filter of nine shared
use-practice terms (e.g., \textit{use case}, \textit{how do you use},
\textit{helps me}) to retain only records describing actual usage
behavior---not model discussion in the abstract.
This two-stage strategy required only \textbf{24 API queries} rather than the
$12 \times 9 = 108$ a cross-product design would need, while preserving the
same logical target (\textit{model-context} AND \textit{use-practice}).

After deduplication by Reddit ID the corpus contained
\textbf{3,292 records}: 1,149 from Stream~1 and 2,143 from Stream~2.

% -----------------------------------------------------------------------------
\subsection{LLM Annotation}
% -----------------------------------------------------------------------------

Each record was submitted to GPT-4o via the OpenAI Batch API (50\% lower cost
than synchronous calls; an incremental resume mechanism handled interrupted
batches).
The annotation prompt asked the model to identify all concrete daily tasks
mentioned in the text and to assign each to one of four themes---
\textit{Creation}, \textit{Information Search}, \textit{Advice},
\textit{Automation}---and ten sub-themes, following the enterprise-use taxonomy
of Brachman et al.~\cite{brachman2024enterprise}.
Per record the schema returned task descriptions, theme and sub-theme
assignments, and evidence spans; multi-task records stored fields
pipe-separated.

% -----------------------------------------------------------------------------
\subsection{Canonical Grouping and Overlap Analysis}
% -----------------------------------------------------------------------------

Free-text task labels were normalized in two passes.
First, 35 hand-crafted regex patterns collapsed surface variants into shared
labels (e.g., \texttt{summar(iz|is)(e|ing|ation)} $\to$
\textit{summarization}).
Second, sentence-transformer embeddings and agglomerative clustering were
applied to the 206 manually reviewed records to merge semantically similar
labels into canonical groups.

This produced \textbf{29 canonical groups} from 39 unique raw labels.
A task group was classified as \textit{shared} if at least one annotated
record per stream mapped to it; \textit{SLM-only} or \textit{LLM-only}
otherwise.
Of the 29 groups, \textbf{19 (65\%) were shared}, 8 were SLM-only, and 2
were LLM-only.
Coverage was 104/104 SLM records and 100/102 LLM records, confirming that
the shared task space accounts for the large majority of observed discourse.

Table~\ref{tab:overlap} lists the ten highest-frequency shared groups.
The high overlap---despite the communities using different model classes in
different infrastructure contexts---supports the premise that everyday task
needs are largely model-agnostic; the \textit{capability} to satisfy those
needs, not the tasks themselves, is what differs.

\begin{table}[t]
  \caption{Top shared canonical task groups ranked by combined mention count.}
  \label{tab:overlap}
  \small
  \begin{tabular}{lrrr}
    \toprule
    Canonical task & SLM & LLM & Total \\
    \midrule
    Creative writing      & 18 & 21 & 39 \\
    Summarization         & 18 & 11 & 29 \\
    Scripting             &  8 & 12 & 20 \\
    Task automation       &  8 & 11 & 19 \\
    Brainstorming         &  4 & 11 & 15 \\
    Transcription         & 11 &  2 & 13 \\
    Translation           &  5 &  3 &  8 \\
    Analyze data          &  2 &  6 &  8 \\
    Planning / scheduling &  2 &  4 &  6 \\
    Agent workflow        &  3 &  3 &  6 \\
    \bottomrule
  \end{tabular}
\end{table}

% -----------------------------------------------------------------------------
\subsection{Asymmetric Patterns and Design Implications}
% -----------------------------------------------------------------------------

Two asymmetries are notable.
\textit{Transcription} is SLM-dominant (11 vs.\ 2 mentions), consistent with
privacy concerns driving local audio processing---an SLM-specific use case
that warrants dedicated scaffolding support not required in the LLM condition.
\textit{Brainstorming} is LLM-dominant (4 vs.\ 11), suggesting that
open-ended ideation is currently under-served by local deployments; the
prompt-variant generator in DG3 directly targets this gap.
The 8 SLM-only groups (\textit{email management}, \textit{RAG},
\textit{knowledge base}, etc.) confirm that privacy-sensitive and
infrastructure-heavy workflows are a distinctive SLM use case beyond the
shared core.

% -----------------------------------------------------------------------------
\subsection{Final Task Shortlist}
% -----------------------------------------------------------------------------

For the comparative study we selected the highest-frequency shared task from
each of the four annotation themes, merging \textit{Advice} and
\textit{Automation} (low-count overlap) into a single planning condition:

\begin{itemize}
  \item \textbf{Task~1 -- Writing} (\textit{Creation} theme):
    \textit{creative writing} --- 18 SLM, 21 LLM mentions.
    Participants produce a practical short piece (email, social post, or
    short description) until the output is something they would ``realistically
    use or send.''
  \item \textbf{Task~2 -- Summarizing} (\textit{Information Search} theme):
    \textit{summarization} --- 18 SLM, 11 LLM mentions.
    Participants summarize a 250-word passage until the result is clear enough
    to ``use instead of re-reading the original.''
  \item \textbf{Task~3 -- Planning} (\textit{Advice} theme):
    \textit{planning / scheduling} --- 2 SLM, 4 LLM mentions.
    Participants produce a day-by-day plan for a given scenario until it feels
    ``realistic, specific, and easy to follow.''
\end{itemize}

The three tasks span a spectrum from open-ended creative generation (highest
LLM advantage expected) to structured output (where constraint scaffolding
may close the gap), enabling a graded test of the ``good enough'' boundary
(RQ2).
